\section{Entornos de trabajo}

El desarrollo del presente proyecto requiere de un entorno colaborativo, flexible y bien estructurado que facilite tanto la construcción como la validación de los distintos componentes del sistema. Para ello, se han seleccionado diversas plataformas que permiten organizar tareas, escribir y probar código, almacenar datos, así como documentar y presentar avances de manera clara y profesional.

\subsection{Plataformas de desarrollo}

 Con el fin de cubrir las distintas etapas del proyecto, desde la implementación hasta la comunicación de resultados, se utilizaron múltiples herramientas que favorecen la eficiencia y el trabajo en equipo. Las siguientes plataformas han sido clave para mantener un flujo de trabajo continuo y bien coordinado.

\begin{itemize}	
	\item\textbf{Visual Studio Code} Comúnmente conocido como VSCode, es un editor de código gratuito, de código abierto y compatible con múltiples sistemas operativos. Fue creado por Microsoft, empresa con amplia trayectoria en el desarrollo de entornos integrados (IDE), lo que se refleja en la eficiencia y funcionalidad que ofrece esta herramienta, muy valorada por la comunidad de desarrolladores \cite{cimas2024vscode}.
	
	Aunque su nombre puede sugerir una relación directa con Visual Studio, en realidad VSCode es una aplicación independiente, construida desde cero con tecnologías diferentes. Mientras Visual Studio se orienta principalmente a tecnologías propias de Microsoft, VSCode destaca por su flexibilidad y capacidad de adaptarse a prácticamente cualquier lenguaje de programación \cite{garcia2024vscode}.
	
	\item\textbf{Github} Es una plataforma basada en la nube donde puedes almacenar, compartir, trabajar y versionar junto con otros usuarios para escribir código \cite{github2025about}. En adición a lo anterior, también permite realizar lo siguiente:
	
	\begin{itemize}
		\item Presentar o compartir el trabajo.
		\item Seguir y administrar los cambios en el código a lo largo del tiempo.
		\item Dejar que otros usuarios revisen el código y realicen sugerencias para mejorarlo.
		\item Colaborar en un proyecto compartido, sin preocuparse de que los cambios afectarán al trabajo de los colaboradores antes de que esté listo para integrarlos.
	\end{itemize}
	
	El versionamiento, una de las características fundamentales de GitHub, es posible gracias al software de código abierto Git, en el que se basa GitHub \cite{github2025about}.
	
	\item\textbf{Canva}	Lanzada en 2013, Canva es una herramienta de diseño y publicación online que tiene como misión poner el poder del diseño al alcance de todo el mundo, para que cualquier persona pueda diseñar lo que quiera y publicar donde quiera \cite{canva2025about}.
		
	\item\textbf{TexStudio}	TeXstudio es un entorno de escritura integrado para la creación de documentos LaTeX. Cuenta con numerosas características como resaltado de sintaxis, visor integrado, comprobación de referencias y varios asistentes. TeXstudio es de código abierto y está disponible para todos los principales sistemas operativos \cite{texstudio}.
	
	
\end{itemize}
